% status: ZERO
\chapter{Numerics: Floating Point and Quantization Formats}\label{app:numerics}
\begin{ChapterIntent}
Ask which values a format represents, how arithmetic rounds, and which
equivalence the application requires. These questions precede bit counts.
Chapters 10--12 provide the detailed low-precision formats and their
measured trade-offs; this appendix establishes the numerical vocabulary.
\end{ChapterIntent}

\section{Finite storage changes arithmetic}
A normal binary floating-point value has the schematic form
\[
 x=(-1)^s\,2^e(1+f),
\]
where $s$ is its sign, $e$ its unbiased exponent and $f$ the stored fraction
interpreted in $[0,1)$. This does not describe zeros, subnormals, infinities
or NaNs. Their encodings are format-specific.

Range and precision solve different problems. A larger range can avoid
overflow without resolving small changes near a large value. Flushing
subnormals matters near zero. Input, multiplication, accumulation and
output formats can differ within one operator; an FP16 label alone leaves
those questions unanswered.

\section{Stable softmax without overflow}
For scores $[1000,999]$, direct exponentiation can overflow. Subtract the
largest score:
\[
 \frac{[e^{1000},e^{999}]}{e^{1000}+e^{999}}
 =\frac{[1,e^{-1}]}{1+e^{-1}}
 \approx[0.731059,0.268941].
\]
Equality here is a real-arithmetic identity; the decimal values are rounded.
Adding the same finite constant to every score does not change exact
probabilities. A masked position contributes no mass. An entirely masked
row has no distribution under this normalization and needs an explicit
policy; blindly evaluating its maximum is not a policy.

A row summing to one can still be wrong. Compare its entries and downstream
weighted output, not merely normalization. Chapter~\ref{ch:flashattention}
uses this contract while changing the evaluation order.

\section{Specify the rounding rule}
Consider illustrative symmetric quantization with scale $s_q=0.5$:
\[
 q=\operatorname{clip}(\operatorname{round}(x/s_q),-7,7),
 \qquad \widehat{x}=s_q q.
\]
For $x=[-1.1,0.2,1.6]$, nearest rounding gives codes $[-2,0,3]$
and reconstruction $[-1,0,1.5]$. Absolute errors are $[0.1,0.2,0.1]$.
Without clipping, nearest rounding bounds error by $s_q/2$.
Clipping removes that bound: $x=10$ reconstructs as $3.5$, with error $6.5$.

At $x=1.25$ the quotient is $2.5$. Ties-to-even gives code 2;
ties-away-from-zero gives 3. The word round is therefore insufficient to
reproduce a conversion. The executable example uses ties-to-even and
tests this boundary. A production encoding must specify saturation,
nonfinite handling and scale computation too.

\section{Metadata is storage}
A hypothetical group of 32 four-bit codes and one 16-bit scale occupies
$32\cdot4+16=144$ bits: 18 bytes, or 4.5 bits per value.
This is an illustrative encoding, not every format named four-bit.
Zero points, secondary scales, padding and outlier storage change the count.
Smaller groups can improve local scaling while increasing metadata cost.

There is a computational bill as well: load and decode metadata, consume
packed codes, and accumulate products. Reduced payload does not guarantee
a faster unsupported kernel. Separate storage savings from measured runtime.

\section{Choose the equivalence}
Floating-point addition is not associative. In an illustrative decimal
system retaining three significant digits, adding $1000$, $1$, $-1000$
left-to-right loses the $1$; cancelling the large terms first retains it.
Binary formats exhibit the same phenomenon at their own spacing boundaries.

A numerical test can require
\[
 |\widehat y-y|\leq\epsilon_{\rm abs}+\epsilon_{\rm rel}|y|.
\]
Declare both tolerances; the absolute term matters near zero. If two logits
differ by a margin $\delta$ and each can move by at most $\epsilon$, their
ordering is guaranteed only when $\delta>2\epsilon$. Small tensor error,
the same greedy token and task quality are thus different tests.

A seed controls random draws, not rounding or reduction order. Quantization
also changes the model's numerical function. Report the intended contract
before deciding that a speedup preserves correctness; Chapter
\ref{ch:fast-wrong-answers} develops this hierarchy.
